\exercise{4.5}{When a supremum is a maximum}

\goal Let $S$ be a nonempty subset of $\R$ that is bounded above. Prove
that if $\sup S\in S$, then $\sup S=\max S$.

\step{Step 1: recall the two requirements for a maximum.}
A number $M$ is the maximum of $S$ exactly when
\begin{enumerate}[label=(\roman*)]
 \item $M$ is an element of $S$;
 \item every $s\in S$ satisfies $s\le M$.
\end{enumerate}
The first requirement concerns membership; the second concerns order.

\step{Step 2: check those requirements for the supremum.}
Since $S$ is nonempty and bounded above, completeness supplies the real
number $M=\sup S$. The hypothesis of the exercise is precisely $M\in S$,
so requirement (i) holds. A supremum is an upper bound, so $s\le M$
for every $s\in S$; requirement (ii) holds as well.

\answer{The number $\sup S$ belongs to $S$ and is at least every member
of $S$. By the definition of maximum,
\(\boxed{\sup S=\max S}\). The additional membership hypothesis is
essential: an upper bound need not be an element of its set.}

\exercise{4.6}{Comparing the infimum and supremum}

\goal For a nonempty bounded set $S\subseteq\R$, prove (a)
$\inf S\le\sup S$, and (b) describe $S$ when these two numbers are equal.

\step{Part (a), Step 1: identify the real bounds and choose an element.}
Because $S$ is nonempty and bounded, the completeness axiom and its
lower-bound version give real numbers
\[
 m=\inf S,\qquad M=\sup S.
\]
Nonemptiness permits us to choose a specific $s_0\in S$.

\step{Part (a), Step 2: apply the meanings of lower and upper bound.}
Since $m$ is a lower bound, $m\le s_0$. Since $M$ is an upper bound,
$s_0\le M$. Transitivity of $\le$ gives
\[
 m\le s_0\le M\quad\Longrightarrow\quad m\le M.
\]

\answer{For every nonempty bounded $S\subseteq\R$,
\(\boxed{\inf S\le\sup S}\). Choosing $s_0$ is why nonemptiness
matters in this argument.}

\step{Part (b), Step 1: let the common value be $c$.}
Assume $\inf S=\sup S=c$. For an arbitrary $s\in S$, the lower-bound
and upper-bound properties give
\[
 c\le s\le c.
\]
By antisymmetry of the order, $s=c$. Since $s$ was arbitrary, every
element of $S$ equals $c$, and hence $S\subseteq\{c\}$.

\step{Part (b), Step 2: use nonemptiness to get the reverse inclusion.}
There is an $s_0\in S$. The preceding step gives $s_0=c$, so $c\in S$.
Thus $\{c\}\subseteq S$, and the two inclusions give $S=\{c\}$.
Conversely, for the singleton $\{c\}$, the element $c$ is both its
maximum and its minimum. Every upper bound must be at least $c$, and
every lower bound must be at most $c$. Therefore both its supremum
and its infimum equal $c$.

\answer{Equality of the infimum and supremum occurs exactly when $S$
is a singleton:
\(\boxed{\inf S=\sup S=c\ \Longleftrightarrow\ S=\{c\}}\).}

\exercise{4.7}{Bounds of a subset and the supremum of a union}

\goal Let $S,T$ be nonempty bounded subsets of $\R$. Prove (a) that
$S\subseteq T$ implies
$\inf T\le\inf S\le\sup S\le\sup T$, and (b) that
$\sup(S\cup T)=\max\{\sup S,\sup T\}$ without assuming either set
contains the other.

\step{Part (a), Step 1: prove the comparison between the infima.}
Assume $S\subseteq T$. Since $\inf T$ is a lower bound of $T$, it
satisfies $\inf T\le t$ for every $t\in T$. Every $s\in S$ is also
an element of $T$, so
\[
 \inf T\le s\qquad\text{for every }s\in S.
\]
Consequently $\inf T$ is a lower bound of $S$. The number $\inf S$
is the \emph{greatest} lower bound of $S$, so every lower bound of
$S$ is at most $\inf S$. In particular,
\[
 \inf T\le\inf S.
\]

\step{Part (a), Step 2: prove the middle comparison.}
Exercise~4.6(a), applied to the nonempty bounded set $S$, gives
\[
 \inf S\le\sup S.
\]

\step{Part (a), Step 3: prove the comparison between the suprema.}
Because $\sup T$ is an upper bound of $T$, every $s\in S\subseteq T$
satisfies $s\le\sup T$. Thus $\sup T$ is an upper bound of $S$.
Since $\sup S$ is the \emph{least} upper bound of $S$,
\[
 \sup S\le\sup T.
\]

\answer{Combining the three inequalities gives
\(\boxed{\inf T\le\inf S\le\sup S\le\sup T}\) whenever
$S\subseteq T$.}

\step{Part (b), Step 1: verify that the union has real bounds.}
Write
\[
 U=S\cup T,\qquad
 L=\min\{\inf S,\inf T\},\qquad
 M=\max\{\sup S,\sup T\}.
\]
The numbers $L,M$ are real. Since $S$ is nonempty, the union $U$ is
nonempty. If $u\in U$, then $u\in S$ or $u\in T$.
In the first case,
\[
 L\le\inf S\le u\le\sup S\le M;
\]
in the second case,
\[
 L\le\inf T\le u\le\sup T\le M.
\]
Therefore $U$ is bounded below by $L$ and above by $M$. In particular,
$\sup U$ exists as a real number.

\step{Part (b), Step 2: prove that $\sup U\le M$.}
Step~1 showed $u\le M$ for every $u\in U$, so $M$ is an upper bound
of $U$. The least-upper-bound property gives $\sup U\le M$.

\step{Part (b), Step 3: prove the inequality in the other direction.}
Since $S\subseteq U$, every member of $S$ is at most $\sup U$.
Hence $\sup U$ is an upper bound of $S$, and $\sup S\le\sup U$.
The same reasoning with $T\subseteq U$ gives $\sup T\le\sup U$.
The maximum $M$ is one of these two numbers; both are at most
$\sup U$. Thus $M\le\sup U$.

\answer{The two comparisons $\sup U\le M$ and $M\le\sup U$ give
\[
 \boxed{\sup(S\cup T)=\max\{\sup S,\sup T\}}.
\]
We used the containments $S\subseteq S\cup T$ and $T\subseteq S\cup T$,
not a containment between $S$ and $T$.}

\exercise{4.8}{Two sets ordered entirely below and above each other}

\goal Let $S,T\subseteq\R$ be nonempty and suppose
\[
 s\le t\qquad\text{for every }s\in S\text{ and every }t\in T.
\]
Show (a) that $S$ is bounded above and $T$ is bounded below, (b) that
$\sup S\le\inf T$, and supply the examples requested in (c) and (d).

\step{Part (a), Step 1: choose one member from each set.}
Because both sets are nonempty, there are $s_0\in S$ and $t_0\in T$.
Keep these two particular numbers fixed for this step.
The hypothesis, with $t=t_0$, says
\[
 s\le t_0\qquad\text{for every }s\in S.
\]
Thus the real number $t_0$ is an upper bound of $S$.
The hypothesis, with $s=s_0$, says
\[
 s_0\le t\qquad\text{for every }t\in T.
\]
Thus the real number $s_0$ is a lower bound of $T$.

\answer{$S$ is bounded above, and $T$ is bounded below. No lower
bound for $S$, or upper bound for $T$, is required by this conclusion.}

\step{Part (b), Step 1: establish existence of the two numbers.}
The set $S$ is nonempty and bounded above, so the completeness axiom
gives $\sup S\in\R$. The set $T$ is nonempty and bounded below, so
Corollary~4.5 gives $\inf T\in\R$.

\step{Part (b), Step 2: turn the supremum into a lower bound of $T$.}
Take an arbitrary $t\in T$. The hypothesis says that $s\le t$ for
every $s\in S$, so $t$ is an upper bound of $S$. Since $\sup S$ is
the least upper bound of $S$, we have $\sup S\le t$.
This is true for every $t\in T$. Therefore $\sup S$ is a lower bound
of $T$. Its greatest lower bound is $\inf T$, so
\[
 \sup S\le\inf T.
\]

\answer{The separation of all pairs of set elements implies
\(\boxed{\sup S\le\inf T}\). The proof does not assume that either
extremal bound belongs to its set.}

\step{Part (c), Step 1: give sets that share an element.}
Take
\[
 S=[0,1],\qquad T=[1,2].
\]
Both sets are nonempty because $1$ belongs to each. If $s\in S$ and
$t\in T$, the interval definitions give $s\le1$ and $1\le t$;
transitivity gives $s\le t$, as required.

\step{Part (c), Step 2: identify their intersection and bounds.}
If $x\in S\cap T$, membership in $S$ gives $x\le1$, and membership
in $T$ gives $1\le x$. Antisymmetry forces $x=1$. Conversely, $1$
belongs to both intervals, so
\[
 S\cap T=\{1\}.
\]
The number $1$ is the maximum of $S$ and the minimum of $T$.
Hence $\sup S=1=\inf T$.

\answer{The example \(\boxed{S=[0,1],\ T=[1,2]}\) meets the pairwise
order condition and has nonempty intersection $\{1\}$.}

\step{Part (d), Step 1: give disjoint sets with a common extremal bound.}
Take
\[
 S=(0,1),\qquad T=(1,2).
\]
The first contains $1/2$, and the second contains $3/2$, so both are
nonempty. For $s\in S$ and $t\in T$, we have $s<1<t$, and hence
$s\le t$. If $x$ belonged to their intersection, we would have both
$x<1$ and $x>1$, which is impossible. Thus $S\cap T=\varnothing$.

\step{Part (d), Step 2: prove directly that $\sup S=1$.}
Every $s\in(0,1)$ satisfies $s<1$, so $1$ is an upper bound of $S$.
To prove it is the least, take any real $u<1$ and put
\[
 v=\max\{0,u\},\qquad s=\frac{v+1}{2}.
\]
Both $0$ and $u$ are less than $1$, so $v<1$. Since $v\ge0$,
the inequalities $v<s<1$ imply $0<s<1$. For completeness,
\[
 s-v=\frac{1-v}{2}>0,\qquad 1-s=\frac{1-v}{2}>0.
\]
Thus $s\in S$ and $s>v\ge u$. Consequently $u$ cannot be an upper
bound of $S$. Every number less than $1$ fails to be an upper bound,
so $\sup S=1$.

\step{Part (d), Step 3: prove directly that $\inf T=1$.}
Every $t\in(1,2)$ satisfies $t>1$, so $1$ is a lower bound of $T$.
For an arbitrary real $\ell>1$, put
\[
 w=\min\{2,\ell\},\qquad t=\frac{1+w}{2}.
\]
Both $2$ and $\ell$ exceed $1$, so $w>1$. We have
\[
 t-1=\frac{w-1}{2}>0,\qquad w-t=\frac{w-1}{2}>0.
\]
Therefore $1<t<w\le2$, which gives $t\in T$. Also $t<w\le\ell$,
so $t<\ell$. Thus $\ell$ cannot be a lower bound of $T$.
Every number greater than $1$ fails to be a lower bound, so $\inf T=1$.

\answer{The example \(\boxed{S=(0,1),\ T=(1,2)}\) has
$\boxed{\sup S=\inf T=1}$ and $\boxed{S\cap T=\varnothing}$.
Equality of the bounds does not force either set to contain the common
bound.}

\exercise{4.9}{Obtain an infimum by negating a supremum}

\goal Complete the proof of Corollary~4.5. For a nonempty set
$S\subseteq\R$ that is bounded below, put
\[
 -S=\{-s:s\in S\},\qquad s_0=\sup(-S).
\]
Prove the two assertions left open in the text:
\begin{align*}
 \text{(1)}\quad&-s_0\le s &&\text{for every }s\in S,\\
 \text{(2)}\quad&t\le s\text{ for every }s\in S
       &&\Longrightarrow\quad t\le-s_0.
\end{align*}

\step{Step 1: justify that $s_0$ exists as a real number.}
Since $S$ is nonempty, choose $s_*\in S$. Then $-s_*\in-S$, so
$-S$ is also nonempty. Since $S$ is bounded below, there is a real
number $m$ with $m\le s$ for every $s\in S$.
Negating this inequality reverses it, so $-s\le-m$ for every $s\in S$.
Every element of $-S$ has the form $-s$ for such an $s$. Therefore
$-m$ is an upper bound of $-S$. The completeness axiom applies and
gives the real number $s_0=\sup(-S)$.

\step{Step 2: prove assertion (1), the lower-bound property.}
Let $s\in S$ be arbitrary. By the definition of $-S$, its negative
$-s$ belongs to $-S$. Because $s_0$ is an upper bound of that set,
\[
 -s\le s_0.
\]
Negate both sides. The direction reverses, and $-(-s)=s$, giving
\[
 -s_0\le s.
\]
This holds for every $s\in S$, so $-s_0$ is a lower bound of $S$.

\step{Step 3: prove assertion (2), the greatest-lower-bound property.}
Let $t$ be any real lower bound of $S$. This means
$t\le s$ for every $s\in S$. Negating each of those inequalities gives
\[
 -s\le-t\qquad\text{for every }s\in S.
\]
Hence $-t$ is an upper bound of $-S$. The least-upper-bound property
of $s_0=\sup(-S)$ now gives
\[
 s_0\le-t.
\]
Negating once more yields $t\le-s_0$. Thus every lower bound $t$ of
$S$ is at most $-s_0$.

\answer{Assertion (1) says $-s_0$ is a lower bound, and assertion (2)
says it is greater than or equal to every lower bound. Together they
prove
\(\boxed{\inf S=-s_0=-\sup(-S)}\). This establishes the lower-bound
form of completeness from its upper-bound form.}

\exercise{4.10}{One natural number meeting two inequalities}

\goal Given a real number $a>0$, find one $n\in\N$ such that
$1/n<a<n$.

\idea It is not enough to find one integer for $1/n<a$ and a different
integer for $a<n$. We choose a single integer larger than both $a$ and
$1/a$.

\step{Step 1: form one positive real target.}
Since $a>0$, its inverse $1/a$ exists and is positive. Define
\[
 B=\max\{a,1/a\}.
\]
The number $B$ is positive because both entries in this maximum are
positive.

\step{Step 2: apply the Archimedean property.}
Use Theorem~4.6 with the positive increment $1$ and the positive target
$B$. It gives a natural number $n$ satisfying
\[
 n\cdot1>B.
\]
Thus $n>B\ge a$, so $a<n$, and $n>B\ge1/a$, so $n>1/a$.

\step{Step 3: convert $n>1/a$ into the reciprocal inequality.}
Multiply $n>1/a$ by the positive number $a$ to get
$na>1$. Because $n\in\N$, we have $n>0$, so division by $n$
preserves strict inequality:
\[
 na>1\quad\Longrightarrow\quad a>\frac1n.
\]

\answer{The same natural number $n$ satisfies both comparisons:
\(\boxed{\frac1n<a<n}\). Both positivity checks are needed when
multiplying and dividing inequalities.}

\exercise{4.11}{Every nonempty open interval contains infinitely many rationals}

\goal If $a,b\in\R$ and $a<b$, use the denseness of $\Q$ to prove
that there are infinitely many rational numbers $r$ with $a<r<b$.

\step{Step 1: define the set whose size is at issue.}
Let
\[
 E=\{r\in\Q:a<r<b\}.
\]
Theorem~4.7, applied to $a<b$, gives at least one rational in this
interval. Thus $E$ is nonempty. This first application alone does not
prove infinitude; we must rule out the possibility that the set is finite.

\step{Step 2: suppose there were only finitely many.}
Assume, for a contradiction, that
\[
 E=\{r_1,r_2,\ldots,r_k\}
\]
for some positive integer $k$. A nonempty finite set of real numbers
has a maximum: compare its first two entries, keep the larger, and
continue through the finitely many entries. Let
\[
 r_* = \max\{r_1,\ldots,r_k\}.
\]
Since $r_*\in E$, it satisfies $a<r_*<b$.

\step{Step 3: use density in a smaller interval.}
The strict inequality $r_*<b$ permits another application of
Theorem~4.7. There is a rational $q$ with
\[
 r_*<q<b.
\]
Because $a<r_*<q$, transitivity gives $a<q$. Together with $q<b$ and
$q\in\Q$, this shows $q\in E$.
But $q>r_*$, contradicting the choice of $r_*$ as the largest member of
$E$.

\answer{The set $E$ is nonempty and cannot be finite. Therefore there
are \textbf{infinitely many rational numbers between $a$ and $b$}.
The same density theorem can always be applied above the largest
member of any proposed finite list.}

\exercise{4.12}{Irrational numbers are dense as well}

\goal Let $I=\R\setminus\Q$. Prove that for any real numbers $a<b$
there is an $x\in I$ such that $a<x<b$.

\step{Step 1: prove that a rational translate of $\sqrt2$ is irrational.}
The irrationality of $\sqrt2$ was established in Section~2.
Let $r\in\Q$ be arbitrary and suppose, for a contradiction, that
$r+\sqrt2$ were rational. Write
\[
 r=\frac mn,\qquad r+\sqrt2=\frac hk,
\]
where $m,n,h,k\in\Z$, $n\ne0$, and $k\ne0$.
Subtracting gives
\[
 \sqrt2=\frac hk-\frac mn=\frac{hn-mk}{kn}.
\]
The numerator $hn-mk$ is an integer, the denominator $kn$ is an
integer, and $kn\ne0$ because both $k$ and $n$ are nonzero.
The last quotient would make $\sqrt2$ rational, a contradiction.
Thus $r+\sqrt2$ is irrational for every $r\in\Q$, proving
\[
 \{r+\sqrt2:r\in\Q\}\subseteq I.
\]

\step{Step 2: shift the desired interval before choosing a rational.}
Adding $-\sqrt2$ to $a<b$ preserves strict inequality:
\[
 a-\sqrt2<b-\sqrt2.
\]
By Theorem~4.7, there is a rational number $r$ satisfying
\[
 a-\sqrt2<r<b-\sqrt2.
\]
The endpoints here are real numbers; they need not themselves be
rational for the density theorem to apply.

\step{Step 3: shift that rational back.}
Define $x=r+\sqrt2$. Step~1 gives $x\in I$. Add $\sqrt2$ to all three
expressions in the preceding two inequalities:
\[
 (a-\sqrt2)+\sqrt2<r+\sqrt2<(b-\sqrt2)+\sqrt2.
\]
Associativity and additive cancellation simplify the endpoints, giving
$a<x<b$.

\answer{The irrational number \(\boxed{x=r+\sqrt2}\), where
$a-\sqrt2<r<b-\sqrt2$ and $r\in\Q$, lies strictly between $a$ and
$b$. The initial shift is what ensures that adding $\sqrt2$ does not
move our chosen point outside the desired interval.}

\exercise{4.13}{Three descriptions of an open interval}

\goal For real $a,b,c$, prove the equivalence of
\begin{align*}
 \text{(i)}\quad&\abs{a-b}<c,\\
 \text{(ii)}\quad&b-c<a<b+c,\\
 \text{(iii)}\quad&a\in(b-c,b+c).
\end{align*}
We interpret $(u,v)$ as the set $\{x\in\R:u<x<v\}$; this set is
empty when $u\ge v$.

\step{Step 1: prove that (i) implies (ii).}
Assume $\abs{a-b}<c$. Exercise~3.7(a), with bounded number $a-b$ and
bound $c$, gives
\[
 -c<a-b<c.
\]
Add $b$ throughout. Strict inequalities are preserved by addition, so
\[
 -c+b<(a-b)+b<c+b.
\]
Associativity and the additive inverse rule simplify the middle
expression to $a$. Commutativity turns the endpoints into $b-c$ and
$b+c$. We have obtained $b-c<a<b+c$, which is (ii).

\step{Step 2: prove that (ii) implies (i).}
Assume $b-c<a<b+c$. Add $-b$ throughout and simplify to obtain
\[
 -c<a-b<c.
\]
Exercise~3.7(a) now gives $\abs{a-b}<c$. This is (i).

\step{Step 3: prove the equivalence of (ii) and (iii).}
By the definition of the open interval,
\[
 (b-c,b+c)=\{x\in\R:b-c<x<b+c\}.
\]
Membership of the particular number $a$ in this set means exactly
that $b-c<a<b+c$. Thus (ii) implies (iii), and (iii) implies (ii).

\step{Step 4: verify explicitly what happens when $c\le0$.}
If $c\le0$, statement (i) is false because every absolute value is
nonnegative: it cannot satisfy $0\le\abs{a-b}<c\le0$.
Also $c\le0$ implies $c\le-c$, so adding $b$ gives
$b+c\le b-c$. Statement (ii) would imply $b-c<b+c$, contrary to
that inequality. Thus (ii) is false. The interval in (iii) has its left
endpoint at least as large as its right endpoint, so it is empty;
statement (iii) is false too. This includes $c=0$, when both endpoints
equal $b$.

\answer{For $c>0$, the steps above show that the three statements are
equivalent; for $c\le0$, all three are false. Consequently, for every
real $a,b,c$,
\(\boxed{\text{(i)}\Longleftrightarrow\text{(ii)}\Longleftrightarrow\text{(iii)}}\).}

\exercise{4.14}{Supremum and infimum of the sum of two sets}

\goal Let $A,B\subseteq\R$ be nonempty and bounded. Define
\[
 A+B=\{a+b:a\in A,\ b\in B\}.
\]
Prove
\[
 \text{(a)}\quad\sup(A+B)=\sup A+\sup B,
 \qquad
 \text{(b)}\quad\inf(A+B)=\inf A+\inf B.
\]
The notation $A+B$ means the set of all such sums, not a union.

\step{Preliminary step: check nonemptiness, boundedness, and existence.}
Since $A$ and $B$ are nonempty and bounded, their extrema as bounds
exist as real numbers. Write
\[
 m_A=\inf A,\quad M_A=\sup A,\qquad
 m_B=\inf B,\quad M_B=\sup B.
\]
Choose $a_0\in A$ and $b_0\in B$. Then $a_0+b_0\in A+B$, so the
sum set is nonempty. For arbitrary $a\in A$ and $b\in B$, we have
\[
 m_A\le a\le M_A,\qquad m_B\le b\le M_B.
\]
Adding the lower comparisons and then the upper comparisons gives
\[
 m_A+m_B\le a+m_B\le a+b\le M_A+b\le M_A+M_B.
\]
Thus $A+B$ is bounded below by $m_A+m_B$ and above by $M_A+M_B$.
Completeness and its lower-bound corollary supply real numbers
\[
 m=\inf(A+B),\qquad M=\sup(A+B).
\]
No assertion that an infimum or supremum belongs to $A$ or $B$ has
been made.

\step{Part (a), Step 1: prove the first inequality.}
We just showed that every $a+b\in A+B$ satisfies
$a+b\le M_A+M_B$. Hence $M_A+M_B$ is an upper bound of $A+B$.
Its least upper bound is $M$, so
\[
 M\le M_A+M_B.
\]

\step{Part (a), Step 2: explain how to approach each separate supremum.}
Let $\varepsilon>0$. Because $\varepsilon/2>0$,
$M_A-\varepsilon/2<M_A$. The number $M_A-\varepsilon/2$ cannot be
an upper bound of $A$: such an upper bound would be strictly smaller
than the least upper bound $M_A$. Therefore there is an
$a_\varepsilon\in A$ with
\[
 M_A-\frac\varepsilon2<a_\varepsilon\le M_A.
\]
The same argument for $B$ gives a $b_\varepsilon\in B$ with
\[
 M_B-\frac\varepsilon2<b_\varepsilon\le M_B.
\]
These are approximation statements, not claims that $M_A\in A$ or
$M_B\in B$.

\step{Part (a), Step 3: prove the reverse inequality using a gap.}
Suppose, for a contradiction, that $M<M_A+M_B$. Set
\[
 \varepsilon=M_A+M_B-M>0.
\]
Choose $a_\varepsilon,b_\varepsilon$ as in Step~2. Adding their two
strict lower bounds gives
\begin{align*}
 a_\varepsilon+b_\varepsilon
   &>\left(M_A-\frac\varepsilon2\right)
       +\left(M_B-\frac\varepsilon2\right),\\
   &=M_A+M_B-\varepsilon,\\
   &=M.
\end{align*}
But $a_\varepsilon+b_\varepsilon\in A+B$, and $M$ is an upper bound
of this set, so the same sum must be at most $M$. This is a
contradiction. Therefore $M_A+M_B\le M$.

\answer{Both $M\le M_A+M_B$ and $M_A+M_B\le M$ hold. By
antisymmetry,
\[
 \boxed{\sup(A+B)=\sup A+\sup B}.
\]
Splitting the positive gap
into two halves leaves enough room to choose one near-maximal element
from each set.}

\step{Part (b), Step 1: prove the first inequality.}
The preliminary step showed $m_A+m_B\le a+b$ for every $a+b\in A+B$.
Thus $m_A+m_B$ is a lower bound of the sum set. Since $m$ is the
greatest lower bound,
\[
 m_A+m_B\le m.
\]

\step{Part (b), Step 2: explain how to approach the separate infima.}
Let $\varepsilon>0$. The number $m_A+\varepsilon/2$ is strictly larger
than $m_A$. It cannot be a lower bound of $A$, since $m_A$ is its
greatest lower bound. Therefore there is an $a_\varepsilon\in A$ with
\[
 m_A\le a_\varepsilon<m_A+\frac\varepsilon2.
\]
Likewise there is a $b_\varepsilon\in B$ with
\[
 m_B\le b_\varepsilon<m_B+\frac\varepsilon2.
\]
Again no attainment of either infimum is being assumed.

\step{Part (b), Step 3: prove the reverse inequality using a gap.}
Suppose, for a contradiction, that $m>m_A+m_B$. Put
\[
 \varepsilon=m-m_A-m_B>0.
\]
Choose $a_\varepsilon,b_\varepsilon$ from Step~2. Adding their two
strict upper bounds gives
\begin{align*}
 a_\varepsilon+b_\varepsilon
   &<\left(m_A+\frac\varepsilon2\right)
       +\left(m_B+\frac\varepsilon2\right),\\
   &=m_A+m_B+\varepsilon,\\
   &=m.
\end{align*}
This sum belongs to $A+B$. But $m$ is a lower bound, so every member
of that set must be at least $m$. The contradiction rules out
$m>m_A+m_B$, giving $m\le m_A+m_B$.

\answer{Combining the two comparisons gives
\(\boxed{\inf(A+B)=\inf A+\inf B}\). The argument applies whether
or not the extremal bounds of the original sets are attained.}

\exercise{4.15}{Comparison using reciprocal natural-number errors}

\goal Suppose $a,b\in\R$ and
\[
 a\le b+\frac1n\qquad\text{for every }n\in\N.
\]
Prove $a\le b$.

\idea If $a>b$, the positive gap $a-b$ can be made larger than one
of the numbers $1/n$. That particular $n$ contradicts the hypothesis.

\step{Step 1: assume the opposite of the desired conclusion.}
Suppose $a\le b$ is false. Total ordering then gives $a>b$, so
\[
 d=a-b>0.
\]

\step{Step 2: obtain a strict reciprocal estimate from the Archimedean property.}
Apply Theorem~4.6 to the positive numbers $d$ and $1$. There is an
$n\in\N$ with $nd>1$. Since $n>0$, division by $n$ gives
\[
 d>\frac1n.
\]
The strict inequality is needed: an estimate $1/n\le d$ alone would
not exclude equality at the next step.

\step{Step 3: compare with the hypothesis for that particular $n$.}
Substitute $d=a-b$ and add $b$ to get
\[
 \frac1n<a-b\quad\Longrightarrow\quad b+\frac1n<a.
\]
The hypothesis applies to every natural number, so it applies to this
one and gives $a\le b+1/n$. These inequalities cannot both hold:
they would give $a\le b+1/n<a$.

\answer{The contradiction proves \(\boxed{a\le b}\).
Exercise~3.8 tested $a$ against every real number above $b$.
Here only the particular numbers $b+1/n$ are tested; the Archimedean
property guarantees that at least one of them detects any hypothetical
positive gap $a-b$.}

\exercise{4.16}{A real number as the supremum of the rationals below it}

\goal For each $a\in\R$, prove
\[
 \sup\{r\in\Q:r<a\}=a.
\]

\step{Step 1: define the set and prove that it is nonempty.}
Fix an arbitrary real $a$ and let
\[
 S_a=\{r\in\Q:r<a\}.
\]
The inequality $a-1<a$ is strict because $1>0$. Apply the denseness
of $\Q$ to the interval $(a-1,a)$. It gives a rational $r_0$ with
\[
 a-1<r_0<a.
\]
Thus $r_0\in S_a$, proving nonemptiness. This construction works
even for $a\le0$; one must not assume that $0$ belongs to $S_a$.

\step{Step 2: prove that $a$ is an upper bound.}
By definition, every $r\in S_a$ satisfies $r<a$, and therefore
$r\le a$. Hence $a$ is an upper bound of $S_a$.
In particular, $S_a$ is nonempty and bounded above, so its real
supremum exists by completeness.

\step{Step 3: prove that every smaller candidate fails to be an upper bound.}
Take any real $u<a$. Denseness of $\Q$ gives a rational $r$ with
\[
 u<r<a.
\]
Since $r\in\Q$ and $r<a$, this $r$ belongs to $S_a$.
The additional inequality $u<r$ shows that $u$ is not an upper bound
of $S_a$: it is exceeded by the specific element $r\in S_a$.

\step{Step 4: express the least-upper-bound conclusion.}
Step~2 showed that $a$ is an upper bound. Step~3 showed that every
number strictly less than $a$ fails to be an upper bound.
Equivalently, if $U$ is any upper bound, $U<a$ is impossible, and
the real order therefore forces $a\le U$.
Thus $a$ is at most every upper bound while being one itself.

\answer{The two defining obligations for a supremum are satisfied, so
\(\boxed{\sup\{r\in\Q:r<a\}=a}\) for every real $a$.
The number $a$ need not be rational, and it is excluded from the set
even when it is rational; the supremum need not be attained.}
